%===============================================================================
%
%===============================================================================
\section{La big picture !}

   Créer un noyau n'est pas une mince affaire, on ne sait donc pas
trop comment s'y prendre au début ! Je vous rassure, plus tard on ne
sait toujours pas, \ldots


   L’idée de cette section est alors de proposer une vision un peu
générale, en survolant les principaux éléments, mais sans rentrer dans
le détail d'aucun d'entre eux.

%-------------------------------------------------------------------------------
%
%-------------------------------------------------------------------------------
\subsection{Au commencement, il y a le matériel}

   Lorsqu'on envisage de réaliser un système d’exploitation, c'est que
l'on souhaite exploiter un système, non ? Donc la  première chose dont
on a besoin, c'est d'une bonne connaissance de ce système cible.

   La première (et pour le moment unique) cible de ManuX, c'est un PC
à base de microprocesseur Intel i386. Je ne rentrerai pas dans la
description de ce système physique, largement documenté de nos jours,
et je suppose que vous avez donc une assez bonne connaissance de

\begin{itemize}
   \item son architecture générale ;
   \item le fonctionnement de son micro processeur ;
   \item sa gestion de la mémoire, des entrées-sorties, des
     périphériques, \ldots
\end{itemize}

   Nous reparlerons un peu de chacun de ces éléments en temps utiles.

   Grâce à cette connaissance, nous sommes capables d'écrire du code
qui permette d'exploiter notre système cible.
   
%-------------------------------------------------------------------------------
%
%-------------------------------------------------------------------------------
\subsection{Ça pourrait glisser, on met les chaînes}

   Un système d'exploitation, c'est du logiciel. Vous allez me dire
que ``{\em Real progammers use cat > a.out}'' je suppose, mais cette
fois je vous propose qu'on se la joue un peu plus confort !

   Nous allons donc faire le choix d'utiliser un langage de
programmation de haut niveau pour écrire ce logiciel, et je vous
propose (impose ?) le langage C. Comme tout ne peut pas être écrit en
C, il nous faudra également utiliser de l'assembleur.

   Il nous faut donc un compilateur C, un assembleur, et un éditeur de
liens. Si nous avons peur de devoir déboguer notre code, il nous faudra
également des outils pour ça.

   En ce qui concerne ManuX, une chaîne de compilation C fondée sur gcc
(ou clang) plus un assembleur comme nasm feront très bien
   l'affaire. Nous décrirons cela dans la section dédiée aux outils.

   Grâce à ces outils, nous voilà équipés pour compiler le code que
nous venons d'écrire afin qu'il puisse être exécuté sur notre système
cible.

%-------------------------------------------------------------------------------
%
%-------------------------------------------------------------------------------
\subsection{Chargez !}

   Nous voila bien partis, \ldots{} nous avons un système matériel et
nous avons écrit du code (un noyau, aussi minimaliste soit-il), que
nous avons compilé de sorte à l'exécuter sur le système.

   Mais pour cela, il faut ``charger le noyau'', c'est-à-dire faire en
sorte qu'il se trouve quelque part ({\em a priori} dans la mémoire
vive du système) pour que le microprocesseur l'exécute.

   Ce travail, peu gratifiant et assez technique, est très dépendant
de notre cible matérielle.

   En ce qui concerne ManuX, nous utiliserons deux ou trois méthodes
classiques sur un (vieux) PC. Nous décrirons cela plus précisément
dans la section {\em Chargement en mémoire et boot}.

\begin{quotation}
   Grâce à cette nouvelle étape, youpi, nous allons enfin pouvoir
exécuter notre noyau sur le système cible !
\end{quotation}

%-------------------------------------------------------------------------------
%
%-------------------------------------------------------------------------------
\subsection{Hello word !}

   Maintenant
